% Generated by trazabilidad/matriz.py
\multicolumn{7}{@{}l}{\textit{Pruebas de la base integrada}} \\
L01 & El script completo se ejecuta sobre una base vacía \textit{(Script DDL)} & Decisión: el script DDL se genera del modelo integrado & \nm{pruebas/\allowbreak postgres/\allowbreak run\_tests.sh\allowbreak :56} & \textit{E:} Carga sin errores \newline \textit{O:} Carga sin errores & P & — \\
M01 & El script de carga encuentra el archivo SQL \textit{(Comprobación 1 (bug 1))} & MySQL: \nm{cargar\_bna.ps1\allowbreak :11}, \nm{GUIA\_EJECUCION.md\allowbreak :83}, \nm{README.md\allowbreak :53} & \nm{pruebas/\allowbreak postgres/\allowbreak pruebas.sql\allowbreak :16} & \textit{E:} — \newline \textit{O:} N/A — prevenido por diseño: no hay script de carga con rutas; el DDL es un único archivo (\nm{sql/bna\_core.\allowbreak sql}) que se carga directamente, y su carga se prueba en L01. & N/A & — \\
M04 & Un depósito actualiza el saldo de la cuenta \textit{(Comprobación 4 (bug 3))} & \parencites[pp.~16, 17, 25 y~32]{BNA2026Cuentas}{BNA2026Core}[pp.~15 y~28]{BNA2026Tarjetas}; MySQL: \nm{99\allowbreak :1996}, \nm{99\allowbreak :388}, \nm{99\allowbreak :1811} & \nm{pruebas/\allowbreak postgres/\allowbreak pruebas.sql\allowbreak :19} & \textit{E:} saldo 5000.00 -> 5500.00 \newline \textit{O:} saldo 5000.00 -> 5500.00 & P & — \\
M05 & Pagar reduce la deuda del crédito \textit{(Comprobación 5 (bug 4))} & \parencite[pp.~3 y~8]{BNA2026Cobranza}; MySQL: \nm{99\allowbreak :1865}, \nm{99\allowbreak :1244} & \nm{pruebas/\allowbreak postgres/\allowbreak pruebas.sql\allowbreak :25} & \textit{E:} deuda 10300.00 -> 7200.00 \newline \textit{O:} deuda 10300.00 -> 7200.00 & P & — \\
M06 & \nm{fn\_uuid\_v7()} guarda la hora correcta \textit{(Comprobación 6 (bug 2))} & MySQL: \nm{99\allowbreak :1699}, \nm{99\allowbreak :1700} & \nm{pruebas/\allowbreak postgres/\allowbreak pruebas.sql\allowbreak :31} & \textit{E:} — \newline \textit{O:} N/A — no aplica: la función no se porta; los UUID los genera \nm{gen\_random\_uuid()}, nativa de PostgreSQL (versión 4, sin marca de tiempo). & N/A & — \\
M07 & Los UUID v7 se ordenan por fecha de creación \textit{(Comprobación 7 (bug 2))} & MySQL: \nm{99\allowbreak :1700}, \nm{99\allowbreak :18}, \nm{diccionario\_datos.md\allowbreak :7} & \nm{pruebas/\allowbreak postgres/\allowbreak pruebas.sql\allowbreak :33} & \textit{E:} — \newline \textit{O:} N/A — no aplica: por la misma razón que M06; el orden de las bitácoras lo da su identificador IDENTITY y su marca de tiempo. & N/A & — \\
M08 & Un identificador debe ser un UUID válido \textit{(Comprobación 8 (bug 7))} & \parencite{BNA2026Core}; MySQL: \nm{99\allowbreak :289}, \nm{99\allowbreak :6} & \nm{pruebas/\allowbreak postgres/\allowbreak pruebas.sql\allowbreak :36} & \textit{E:} Rechazo: invalid input syntax for type uuid \newline \textit{O:} Rechazo: invalid input syntax for type uuid: \textquotedbl{}c2\textquotedbl{} & P & — \\
M09a & La orden no abona en la cuenta de otro cliente \textit{(Comprobación 9 (bug 7))} & \parencites[pp.~11 y~12]{BNA2026Desembolso}[p.~8]{BNA2026Evaluacion}; MySQL: \nm{99\allowbreak :1106}, \nm{99\allowbreak :1110}; Decisión: BNA-CUENTA-DESTINO & \nm{pruebas/\allowbreak postgres/\allowbreak pruebas.sql\allowbreak :41} & \textit{E:} Rechazo: BNA-CUENTA-DESTINO.*no pertenece \newline \textit{O:} Rechazo: BNA-CUENTA-DESTINO: la cuenta CTA-PEN-1 no pertenece al cliente C002 de la solicitud & P & — \\
M09b & La orden no abona en una cuenta de otra moneda \textit{(Comprobación 9 (bug 7))} & \parencites[pp.~11 y~12]{BNA2026Desembolso}[p.~8]{BNA2026Evaluacion}; MySQL: \nm{99\allowbreak :1106}, \nm{99\allowbreak :1110}; Decisión: BNA-CUENTA-DESTINO & \nm{pruebas/\allowbreak postgres/\allowbreak pruebas.sql\allowbreak :47} & \textit{E:} Rechazo: BNA-CUENTA-DESTINO.*USD \newline \textit{O:} Rechazo: BNA-CUENTA-DESTINO: la cuenta CTA-USD-1 está en USD y la solicitud en PEN & P & — \\
M09c & La orden no supera el límite aprobado \textit{(Comprobación 9 (bug 7))} & \parencites[pp.~11 y~12]{BNA2026Desembolso}[p.~8]{BNA2026Evaluacion}; MySQL: \nm{99\allowbreak :1106}, \nm{99\allowbreak :1110}; Decisión: BNA-IMPORTE; Decisión: el importe de la orden no supera el límite aprobado (supuesto S12) & \nm{pruebas/\allowbreak postgres/\allowbreak pruebas.sql\allowbreak :53} & \textit{E:} Rechazo: BNA-IMPORTE \newline \textit{O:} Rechazo: BNA-IMPORTE: el importe 99999.00 supera el límite aprobado 10000.00 & P & — \\
M10a & Un retiro no deja el saldo en negativo \textit{(Comprobación 10 (bug 7))} & \parencites[pp.~16, 17, 25 y~32]{BNA2026Cuentas}{BNA2026Core}[pp.~15 y~28]{BNA2026Tarjetas}; MySQL: \nm{99\allowbreak :450}, \nm{99\allowbreak :430}, \nm{99\allowbreak :433}; Decisión: BNA-SALDO & \nm{pruebas/\allowbreak postgres/\allowbreak pruebas.sql\allowbreak :60} & \textit{E:} Rechazo: BNA-SALDO \newline \textit{O:} Rechazo: BNA-SALDO: saldo disponible insuficiente en CTA-PEN-2 (5000.00 < 6000.00) & P & — \\
M10b & Un movimiento va en la moneda de su cuenta \textit{(Comprobación 10 (bug 7))} & \parencites[pp.~16, 17, 25 y~32]{BNA2026Cuentas}{BNA2026Core}[pp.~15 y~28]{BNA2026Tarjetas}; MySQL: \nm{99\allowbreak :450}, \nm{99\allowbreak :430}, \nm{99\allowbreak :433}; Decisión: BNA-MONEDA & \nm{pruebas/\allowbreak postgres/\allowbreak pruebas.sql\allowbreak :63} & \textit{E:} Rechazo: BNA-MONEDA \newline \textit{O:} Rechazo: BNA-MONEDA: el asiento está en USD y la cuenta CTA-PEN-2 en PEN & P & — \\
M11 & El saldo disponible no supera al contable \textit{(Comprobación 11 (bug 7))} & MySQL: \nm{99\allowbreak :401}; Decisión: el saldo disponible no supera al contable & \nm{pruebas/\allowbreak postgres/\allowbreak pruebas.sql\allowbreak :67} & \textit{E:} Rechazo: \nm{ck\_cuenta\_ahorro\_1} \newline \textit{O:} Rechazo: new row for relation \textquotedbl{}\nm{cuenta\_ahorro}\textquotedbl{} violates check constraint \textquotedbl{}\nm{ck\_cuenta\_ahorro\_1}\textquotedbl{} & P & — \\
M12 & El total de la cuota es la suma de sus partes \textit{(Comprobación 12 (bug 7))} & \parencite[p.~2]{BNA2026Cobranza}; MySQL: \nm{99\allowbreak :1269}, \nm{99\allowbreak :1277}; Decisión: el total de la cuota se deriva de sus partes (normalización) & \nm{pruebas/\allowbreak postgres/\allowbreak pruebas.sql\allowbreak :71} & \textit{E:} — \newline \textit{O:} N/A — prevenido por diseño: CUOTA no guarda el total; es un dato derivado de capital, interés y otros conceptos, eliminado en la normalización. & N/A & — \\
M13a & No se aplica más de lo pagado \textit{(Comprobación 13 (bug 4))} & \parencite[pp.~3 y~8]{BNA2026Cobranza}; MySQL: \nm{99\allowbreak :1853}, \nm{99\allowbreak :1867}; Decisión: BNA-PAGO; Decisión: una aplicación no supera el monto del pago & \nm{pruebas/\allowbreak postgres/\allowbreak pruebas.sql\allowbreak :74} & \textit{E:} Rechazo: BNA-PAGO.*se aplicarían \newline \textit{O:} Rechazo: BNA-PAGO: se aplicarían 3500.00 de un pago de 3100.00 & P & — \\
M13b & No se aplica más del saldo de la cuota \textit{(Comprobación 13 (bug 4))} & \parencite[pp.~3 y~8]{BNA2026Cobranza}; MySQL: \nm{99\allowbreak :1853}, \nm{99\allowbreak :1867}; Decisión: BNA-PAGO; Decisión: una aplicación no supera el saldo de la cuota & \nm{pruebas/\allowbreak postgres/\allowbreak pruebas.sql\allowbreak :77} & \textit{E:} Rechazo: BNA-PAGO.*solo debe \newline \textit{O:} Rechazo: BNA-PAGO: se aplican 3101.00 a la cuota 1, que solo debe 3100.00 & P & — \\
M14a & Una orden genera un solo crédito \textit{(Comprobación 14 (bug 7))} & \parencites[p.~2]{BNA2026Cobranza}[p.~4]{BNA2026Solicitud}; MySQL: \nm{99\allowbreak :1251}, \nm{99\allowbreak :1238} & \nm{pruebas/\allowbreak postgres/\allowbreak pruebas.sql\allowbreak :81} & \textit{E:} Rechazo: \nm{uk\_credito} \newline \textit{O:} Rechazo: duplicate key value violates unique constraint \textquotedbl{}\nm{uk\_credito\_1}\textquotedbl{} & P & — \\
M14b & Una solicitud tiene una sola orden confirmada, aun en otra ronda \textit{(Comprobación 14 (bug 7))} & \parencites[pp.~11 y~12]{BNA2026Desembolso}[p.~8]{BNA2026Evaluacion}; MySQL: \nm{99\allowbreak :1251}, \nm{99\allowbreak :1238}; Decisión: BNA-ORDEN-UNICA & \nm{pruebas/\allowbreak postgres/\allowbreak pruebas.sql\allowbreak :84} & \textit{E:} Rechazo: BNA-ORDEN-UNICA \newline \textit{O:} Rechazo: BNA-ORDEN-UNICA: la solicitud b0000000-0000-4000-8000-000000000001 ya tiene una orden CONFIRMADA & P & — \\
M15a & Los últimos cuatro dígitos son numéricos \textit{(Comprobación 15 (bug 7))} & \parencite[p.~17]{BNA2026Tarjetas}; MySQL: \nm{99\allowbreak :527}, \nm{99\allowbreak :531} & \nm{pruebas/\allowbreak postgres/\allowbreak pruebas.sql\allowbreak :94} & \textit{E:} Rechazo: \nm{ck\_tarjeta\_2} \newline \textit{O:} Rechazo: new row for relation \textquotedbl{}tarjeta\textquotedbl{} violates check constraint \textquotedbl{}\nm{ck\_tarjeta\_2}\textquotedbl{} & P & — \\
M15b & La tarjeta vence después de emitirse \textit{(Comprobación 15 (bug 7))} & MySQL: \nm{99\allowbreak :527}, \nm{99\allowbreak :531}; Decisión: la tarjeta vence después de emitirse & \nm{pruebas/\allowbreak postgres/\allowbreak pruebas.sql\allowbreak :98} & \textit{E:} Rechazo: \nm{ck\_tarjeta\_1} \newline \textit{O:} Rechazo: new row for relation \textquotedbl{}tarjeta\textquotedbl{} violates check constraint \textquotedbl{}\nm{ck\_tarjeta\_1}\textquotedbl{} & P & — \\
M16 & La matriz no es una credencial de la tarjeta \textit{(Comprobación 16 (bug 7))} & \parencite[p.~18]{BNA2026Tarjetas}; MySQL: \nm{99\allowbreak :563}, \nm{99\allowbreak :587} & \nm{pruebas/\allowbreak postgres/\allowbreak pruebas.sql\allowbreak :103} & \textit{E:} Rechazo: \nm{ck\_credencial\_1} \newline \textit{O:} Rechazo: new row for relation \textquotedbl{}credencial\textquotedbl{} violates check constraint \textquotedbl{}\nm{ck\_credencial\_1}\textquotedbl{} & P & — \\
M17 & La coordenada pedida existe en la matriz \textit{(Comprobación 17 (bug 7))} & \parencite[p.~20]{BNA2026Tarjetas}; MySQL: \nm{99\allowbreak :639}, \nm{99\allowbreak :514} & \nm{pruebas/\allowbreak postgres/\allowbreak pruebas.sql\allowbreak :108} & \textit{E:} Rechazo: \nm{ck\_intento\_autenticacion\_1} \newline \textit{O:} Rechazo: new row for relation \textquotedbl{}\nm{intento\_autenticacion}\textquotedbl{} violates check constraint \textquotedbl{}\nm{ck\_intento\_autenticacion\_1}\textquotedbl{} & P & — \\
M18 & Una referencia apunta a un registro real \textit{(Comprobación 18 (bug 7))} & \parencite[p.~3]{BNA2026Solicitud}; MySQL: \nm{99\allowbreak :766}, \nm{99\allowbreak :881} & \nm{pruebas/\allowbreak postgres/\allowbreak pruebas.sql\allowbreak :113} & \textit{E:} Rechazo: violates foreign key constraint \newline \textit{O:} Rechazo: insert or update on table \textquotedbl{}evaluacion\textquotedbl{} violates foreign key constraint \textquotedbl{}\nm{fk\_evaluacion\_1}\textquotedbl{} & P & — \\
M19 & La auditoría solo registra eventos sobre registros reales \textit{(Comprobación 19)} & MySQL: \nm{99\allowbreak :259}, \nm{99\allowbreak :260}; Decisión: la entidad auditada es un catálogo (en Cobranza, p. 4, la referencia es libre) & \nm{pruebas/\allowbreak postgres/\allowbreak pruebas.sql\allowbreak :118} & \textit{E:} Rechazo: \nm{fk\_evento\_auditoria\_2} \newline \textit{O:} Rechazo: insert or update on table \textquotedbl{}\nm{evento\_auditoria}\textquotedbl{} violates foreign key constraint \textquotedbl{}\nm{fk\_evento\_auditoria\_2}\textquotedbl{} & P & — \\
M20a & Un asiento descuadrado se rechaza al confirmar \textit{(Comprobación 20 (bug 7))} & \parencites{BNA2026Core}[p.~16]{BNA2026Cuentas}; MySQL: \nm{99\allowbreak :250}, \nm{99\allowbreak :224}; Decisión: BNA-PARTIDA-DOBLE & \nm{pruebas/\allowbreak postgres/\allowbreak pruebas.sql\allowbreak :123} & \textit{E:} Rechazo: BNA-PARTIDA-DOBLE \newline \textit{O:} Rechazo: BNA-PARTIDA-DOBLE: el asiento 9a000000-0000-4000-8000-000000000001 no cuadra (debe 100.00, haber 0.00) & P & — \\
M21 & Una reversión la autoriza otra persona \textit{(Comprobación 21)} & MySQL: \nm{99\allowbreak :1201}, \nm{99\allowbreak :1496}, \nm{99\allowbreak :1497} & \nm{pruebas/\allowbreak postgres/\allowbreak pruebas.sql\allowbreak :129} & \textit{E:} — \newline \textit{O:} N/A — prevenido por diseño: el modelo no guarda usuarios ni empleados; el solicitante y la autorización son identificadores que emite Seguridad, que es quien controla a las personas. & N/A & — \\
M22 & Un retiro se cuenta una sola vez en el tope y en el saldo \textit{(Comprobación 22 (bug 5))} & \parencites[pp.~16, 17, 25 y~32]{BNA2026Cuentas}{BNA2026Core}[pp.~15 y~28]{BNA2026Tarjetas}; MySQL: \nm{99\allowbreak :1801}, \nm{99\allowbreak :1816} & \nm{pruebas/\allowbreak postgres/\allowbreak pruebas.sql\allowbreak :132} & \textit{E:} consumido 600.00; contable 4400.00; disponible 4400.00; movimiento vinculado \newline \textit{O:} consumido 600.00; contable 4400.00; disponible 4400.00; movimiento vinculado & P & — \\
M23 & Una autorización aprobada después también respeta el tope \textit{(Comprobación 23 (bug 5))} & \parencite[pp.~15, 21, 23, 26 y~27]{BNA2026Tarjetas}; MySQL: \nm{99\allowbreak :1796}, \nm{99\allowbreak :1799} & \nm{pruebas/\allowbreak postgres/\allowbreak pruebas.sql\allowbreak :153} & \textit{E:} Rechazo: ERR-TAR-007 \newline \textit{O:} Rechazo: ERR-TAR-007: tope diario excedido (RN-TOP-001); disponible hoy 900.00 PEN & P & — \\
M24a & Se respeta el tope diario del cajero (movimiento) \textit{(Comprobación 24 (bug 5))} & \parencites[pp.~16, 17, 25 y~32]{BNA2026Cuentas}{BNA2026Core}[pp.~15 y~28]{BNA2026Tarjetas}; MySQL: \nm{99\allowbreak :1802}, \nm{99\allowbreak :1412} & \nm{pruebas/\allowbreak postgres/\allowbreak pruebas.sql\allowbreak :162} & \textit{E:} Rechazo: RN-TOP-001 \newline \textit{O:} Rechazo: RN-TOP-001: tope diario de ATM excedido en CTA-PEN-1 (consumido 0.00, tope 1500.00, solicitado 1600.00) & P & — \\
M24b & Se respeta el tope diario del cajero (autorización) \textit{(Comprobación 24 (bug 5))} & \parencite[pp.~15, 21, 23, 26 y~27]{BNA2026Tarjetas}; MySQL: \nm{99\allowbreak :1802}, \nm{99\allowbreak :1412} & \nm{pruebas/\allowbreak postgres/\allowbreak pruebas.sql\allowbreak :165} & \textit{E:} Rechazo: ERR-TAR-007 \newline \textit{O:} Rechazo: ERR-TAR-007: tope diario excedido (RN-TOP-001); disponible hoy 1500.00 PEN & P & — \\
M24c & Dos autorizaciones aprobadas sin movimiento no superan juntas el tope \textit{(Comprobación 24 (bug 5))} & \parencite[pp.~15, 21, 23, 26 y~27]{BNA2026Tarjetas}; MySQL: \nm{99\allowbreak :1802}, \nm{99\allowbreak :1412} & \nm{pruebas/\allowbreak postgres/\allowbreak pruebas.sql\allowbreak :170} & \textit{E:} Rechazo: ERR-TAR-007 \newline \textit{O:} Rechazo: ERR-TAR-007: tope diario excedido (RN-TOP-001); disponible hoy 500.00 PEN & P & — \\
M25 & \nm{v\_tope\_restante\_hoy} muestra el tope de cada tarjeta y cuenta \textit{(Comprobación 25 (bug 5))} & \parencite[pp.~23, 24, 27 y~28]{BNA2026Tarjetas}; MySQL: \nm{99\allowbreak :1568}, \nm{99\allowbreak :1569} & \nm{pruebas/\allowbreak postgres/\allowbreak pruebas.sql\allowbreak :177} & \textit{E:} T1 1500.00; T2 900.00 \newline \textit{O:} a1 1500.00; a2 900.00 & P & — \\
M26 & Al bloquear una tarjeta se guarda el motivo \textit{(Comprobación 26 (bug 6))} & \parencite[pp.~18 y~22]{BNA2026Tarjetas}; MySQL: \nm{99\allowbreak :1827}, \nm{99\allowbreak :1832}, \nm{99\allowbreak :1736} & \nm{pruebas/\allowbreak postgres/\allowbreak pruebas.sql\allowbreak :187} & \textit{E:} Robo reportado / CLIENTE \newline \textit{O:} Robo reportado / CLIENTE & P & — \\
M28 & El diagrama muestra todas las relaciones \textit{(Comprobación 28)} & MySQL: \nm{mer\_bna.mermaid\allowbreak :2}, \nm{mer\_bna.mermaid\allowbreak :63} & \nm{pruebas/\allowbreak postgres/\allowbreak pruebas.sql\allowbreak :195} & \textit{E:} — \newline \textit{O:} N/A — prevenido por diseño: los diagramas, las tablas del informe y el DDL se generan del mismo modelo (\nm{modelo/modelo\_bna.\allowbreak py}), de modo que no pueden divergir. & N/A & — \\
F01 & Solo se emite tarjeta a un cliente ACTIVO \textit{(\nm{trg\_tarjeta\_bi} (RF-TAR-005))} & \parencite[pp.~23 y~31]{BNA2026Tarjetas} & \nm{pruebas/\allowbreak postgres/\allowbreak pruebas.sql\allowbreak :201} & \textit{E:} Rechazo: ERR-TAR-001 \newline \textit{O:} Rechazo: ERR-TAR-001: no se emite tarjeta al cliente C004 en estado BLOQUEADO & P & — \\
F02 & La emisión queda en el historial y en la bandeja hacia Seguridad \textit{(\nm{trg\_tarjeta\_ai} (RN-SEG-002))} & \parencite[pp.~22 y~23]{BNA2026Tarjetas} & \nm{pruebas/\allowbreak postgres/\allowbreak pruebas.sql\allowbreak :205} & \textit{E:} historial 1; evento \nm{EMISION\_TARJETA} 1 \newline \textit{O:} historial 1; evento \nm{EMISION\_TARJETA} 1 & P & — \\
F03 & Una tarjeta CANCELADA no cambia de estado \textit{(\nm{trg\_tarjeta\_bu} (RN-TAR-003))} & \parencite[pp.~23 y~29]{BNA2026Tarjetas} & \nm{pruebas/\allowbreak postgres/\allowbreak pruebas.sql\allowbreak :215} & \textit{E:} Rechazo: ERR-TAR-002 \newline \textit{O:} Rechazo: ERR-TAR-002: la tarjeta a0000000-0000-4000-8000-000000000003 está CANCELADA y no cambia de estado & P & — \\
F04 & Una tarjeta VENCIDA solo se cancela \textit{(\nm{trg\_tarjeta\_bu} (RN-TAR-003))} & \parencite[pp.~23 y~29]{BNA2026Tarjetas} & \nm{pruebas/\allowbreak postgres/\allowbreak pruebas.sql\allowbreak :219} & \textit{E:} Rechazo: ERR-TAR-002 \newline \textit{O:} Rechazo: ERR-TAR-002: una tarjeta VENCIDA solo se cancela (se pidió ACTIVA) & P & — \\
F05 & El cliente no puede desbloquear su tarjeta \textit{(\nm{trg\_tarjeta\_bu} (RN-TAR-004))} & \parencite[pp.~23 y~29]{BNA2026Tarjetas} & \nm{pruebas/\allowbreak postgres/\allowbreak pruebas.sql\allowbreak :223} & \textit{E:} Rechazo: ERR-TAR-004 \newline \textit{O:} Rechazo: ERR-TAR-004: el origen CLIENTE no puede desbloquear la tarjeta a0000000-0000-4000-8000-000000000001 & P & — \\
F06 & El desbloqueo por un operador reinicia los intentos \textit{(\nm{trg\_tarjeta\_au} (RN-TAR-004))} & \parencite[pp.~23 y~29]{BNA2026Tarjetas} & \nm{pruebas/\allowbreak postgres/\allowbreak pruebas.sql\allowbreak :227} & \textit{E:} ACTIVA; intentos 0 \newline \textit{O:} ACTIVA; intentos 0 & P & — \\
F07 & El estado de la tarjeta solo cambia registrando el historial \textit{(Historial (RN-SEG-002))} & \parencite[pp.~23 y~29]{BNA2026Tarjetas}; Decisión: el estado de la tarjeta cambia solo por su historial (en Tarjetas, p. 23, el historial se escribe después del cambio) & \nm{pruebas/\allowbreak postgres/\allowbreak pruebas.sql\allowbreak :239} & \textit{E:} Rechazo: RN-SEG-002 \newline \textit{O:} Rechazo: RN-SEG-002: el estado de la tarjeta a0000000-0000-4000-8000-000000000001 solo cambia registrando el historial & P & — \\
F08 & El tercer intento fallido bloquea la tarjeta y avisa a Seguridad \textit{(\nm{trg\_intento\_ai} (RN-TAR-004))} & \parencite[pp.~19 y~23]{BNA2026Tarjetas} & \nm{pruebas/\allowbreak postgres/\allowbreak pruebas.sql\allowbreak :242} & \textit{E:} BLOQUEADA; evento \nm{BLOQUEO\_POR\_INTENTOS} 1 \newline \textit{O:} BLOQUEADA; evento \nm{BLOQUEO\_POR\_INTENTOS} 1 & P & — \\
F09 & Un intento exitoso reinicia el contador \textit{(\nm{trg\_intento\_ai} (RN-TAR-004))} & \parencite[pp.~19 y~23]{BNA2026Tarjetas} & \nm{pruebas/\allowbreak postgres/\allowbreak pruebas.sql\allowbreak :252} & \textit{E:} intentos 0 \newline \textit{O:} intentos 0 & P & — \\
F10 & El tercer fallo de la matriz la bloquea junto con la tarjeta \textit{(\nm{trg\_intento\_ai} (RN-TAR-004))} & \parencite[pp.~19 y~23]{BNA2026Tarjetas} & \nm{pruebas/\allowbreak postgres/\allowbreak pruebas.sql\allowbreak :260} & \textit{E:} matriz BLOQUEADA; tarjeta BLOQUEADA \newline \textit{O:} matriz BLOQUEADA; tarjeta BLOQUEADA & P & — \\
F11a & La bitácora de intentos no se modifica \textit{(\nm{trg\_intento\_bu} (RN-SEG-002))} & \parencites{BNA2026Core}[p.~23]{BNA2026Tarjetas}[p.~13]{BNA2026Desembolso} & \nm{pruebas/\allowbreak postgres/\allowbreak pruebas.sql\allowbreak :269} & \textit{E:} Rechazo: ERR-TAR-009 \newline \textit{O:} Rechazo: ERR-TAR-009: \nm{tarjetas.\allowbreak intento\_autenticacion} es de solo inserción (UPDATE no permitido) & P & — \\
F11b & El historial de estados no se borra \textit{(\nm{trg\_historial\_bd} (RN-SEG-002))} & \parencites{BNA2026Core}[p.~23]{BNA2026Tarjetas}[p.~13]{BNA2026Desembolso} & \nm{pruebas/\allowbreak postgres/\allowbreak pruebas.sql\allowbreak :274} & \textit{E:} Rechazo: ERR-TAR-009 \newline \textit{O:} Rechazo: ERR-TAR-009: \nm{tarjetas.\allowbreak historial\_estado\_tarjeta} es de solo inserción (DELETE no permitido) & P & — \\
F11c & El historial no se vacía con TRUNCATE \textit{(\nm{trg\_historial} (RN-SEG-002))} & \parencites{BNA2026Core}[p.~23]{BNA2026Tarjetas}[p.~13]{BNA2026Desembolso}; Decisión: el historial tampoco se vacía con TRUNCATE & \nm{pruebas/\allowbreak postgres/\allowbreak pruebas.sql\allowbreak :277} & \textit{E:} Rechazo: ERR-TAR-009 \newline \textit{O:} Rechazo: ERR-TAR-009: \nm{tarjetas.\allowbreak historial\_estado\_tarjeta} es de solo inserción (TRUNCATE no permitido) & P & — \\
F11d & La auditoría no se modifica \textit{(Bitácora de auditoría)} & \parencites{BNA2026Core}[p.~23]{BNA2026Tarjetas}[p.~13]{BNA2026Desembolso}; Decisión: BNA-SOLO-INSERCION & \nm{pruebas/\allowbreak postgres/\allowbreak pruebas.sql\allowbreak :280} & \textit{E:} Rechazo: BNA-SOLO-INSERCION \newline \textit{O:} Rechazo: BNA-SOLO-INSERCION: \nm{creditos.\allowbreak evento\_auditoria} es de solo inserción (UPDATE no permitido) & P & — \\
F11e & Un asiento no se modifica \textit{(Asientos inmutables)} & \parencites{BNA2026Core}[p.~23]{BNA2026Tarjetas}[p.~13]{BNA2026Desembolso}; Decisión: BNA-SOLO-INSERCION & \nm{pruebas/\allowbreak postgres/\allowbreak pruebas.sql\allowbreak :283} & \textit{E:} Rechazo: BNA-SOLO-INSERCION \newline \textit{O:} Rechazo: BNA-SOLO-INSERCION: \nm{core.\allowbreak detalle\_asiento} es de solo inserción (UPDATE no permitido) & P & — \\
F11f & El historial de la solicitud no se borra \textit{(Historial de la solicitud)} & \parencites{BNA2026Core}[p.~23]{BNA2026Tarjetas}[p.~13]{BNA2026Desembolso}; Decisión: BNA-SOLO-INSERCION; Decisión: el historial de la solicitud es de solo inserción & \nm{pruebas/\allowbreak postgres/\allowbreak pruebas.sql\allowbreak :286} & \textit{E:} Rechazo: BNA-SOLO-INSERCION \newline \textit{O:} Rechazo: BNA-SOLO-INSERCION: \nm{creditos.\allowbreak historial\_estado\_solicitud} es de solo inserción (DELETE no permitido) & P & — \\
F12a & No se autoriza con una tarjeta bloqueada \textit{(\nm{trg\_autorizacion\_bi} (RN-TAR-003))} & \parencite[pp.~15, 21, 23, 26 y~27]{BNA2026Tarjetas} & \nm{pruebas/\allowbreak postgres/\allowbreak pruebas.sql\allowbreak :290} & \textit{E:} Rechazo: ERR-TAR-003.*no es operativa \newline \textit{O:} Rechazo: ERR-TAR-003: la tarjeta a0000000-0000-4000-8000-000000000001 no es operativa (estado BLOQUEADA) & P & — \\
F12b & No se autoriza con una tarjeta vencida \textit{(\nm{trg\_autorizacion\_bi} (RN-TAR-003))} & \parencite[pp.~15, 21, 23, 26 y~27]{BNA2026Tarjetas} & \nm{pruebas/\allowbreak postgres/\allowbreak pruebas.sql\allowbreak :295} & \textit{E:} Rechazo: ERR-TAR-003.*venció \newline \textit{O:} Rechazo: ERR-TAR-003: la tarjeta a0000000-0000-4000-8000-0000000000f3 venció el 2021-01-31 & P & — \\
F12c & El cajero no opera una cuenta mancomunada conjunta \textit{(\nm{trg\_autorizacion\_bi} (RN-ATM-005))} & \parencite[pp.~15, 21, 23, 26 y~27]{BNA2026Tarjetas} & \nm{pruebas/\allowbreak postgres/\allowbreak pruebas.sql\allowbreak :302} & \textit{E:} Rechazo: ERR-TAR-008 \newline \textit{O:} Rechazo: ERR-TAR-008: la cuenta CTA-CONJ no puede operarse en cajero (RN-ATM-005) & P & — \\
F12d & La Billetera Digital no opera en dólares \textit{(\nm{trg\_autorizacion\_bi} (RN-BD-002))} & \parencite[pp.~15, 21, 23, 26 y~27]{BNA2026Tarjetas} & \nm{pruebas/\allowbreak postgres/\allowbreak pruebas.sql\allowbreak :306} & \textit{E:} Rechazo: ERR-TAR-006 \newline \textit{O:} Rechazo: ERR-TAR-006: el canal BD no opera en USD (RN-BD-002) & P & — \\
F12e & No se autoriza en una terminal inactiva \textit{(\nm{trg\_autorizacion\_bi})} & \parencite[pp.~15, 21, 23, 26 y~27]{BNA2026Tarjetas}; Decisión: solo autoriza una terminal activa & \nm{pruebas/\allowbreak postgres/\allowbreak pruebas.sql\allowbreak :310} & \textit{E:} Rechazo: ERR-TAR-006.*terminal \newline \textit{O:} Rechazo: ERR-TAR-006: la terminal ATM-AQP-OFF no está activa & P & — \\
F12f & No se autoriza más que el saldo disponible \textit{(\nm{sp\_validar\_reglas\_autorizacion})} & \parencite[pp.~15, 21, 23, 26 y~27]{BNA2026Tarjetas}; Decisión: BNA-SALDO & \nm{pruebas/\allowbreak postgres/\allowbreak pruebas.sql\allowbreak :314} & \textit{E:} Rechazo: BNA-SALDO \newline \textit{O:} Rechazo: BNA-SALDO: saldo disponible insuficiente en CTA-PEN-2 (5000.00 < 8000.00) & P & — \\
F12g & Una autorización rechazada no pasa a autorizada \textit{(\nm{trg\_autorizacion\_bu} (RN-ATM-007))} & \parencite[pp.~15, 21, 23, 26 y~27]{BNA2026Tarjetas} & \nm{pruebas/\allowbreak postgres/\allowbreak pruebas.sql\allowbreak :318} & \textit{E:} Rechazo: RN-ATM-007 \newline \textit{O:} Rechazo: RN-ATM-007: transición no permitida de RECHAZADA a AUTORIZADA & P & — \\
F12h & Una autorización válida se registra \textit{(\nm{trg\_autorizacion\_bi} (C06))} & \parencite[pp.~15, 21, 23, 26 y~27]{BNA2026Tarjetas} & \nm{pruebas/\allowbreak postgres/\allowbreak pruebas.sql\allowbreak :323} & \textit{E:} registrada \newline \textit{O:} registrada & P & — \\
F12i & Dos autorizaciones aprobadas sin movimiento no superan juntas el saldo \textit{(\nm{sp\_validar\_reglas\_autorizacion})} & \parencite[pp.~15, 21, 23, 26 y~27]{BNA2026Tarjetas}; Decisión: BNA-SALDO & \nm{pruebas/\allowbreak postgres/\allowbreak pruebas.sql\allowbreak :330} & \textit{E:} Rechazo: BNA-SALDO \newline \textit{O:} Rechazo: BNA-SALDO: saldo disponible insuficiente en CTA-PEN-2 (2000.00 < 3000.00) & P & — \\
F12j & Extornar una autorización sin movimiento libera el tope y el saldo \textit{(\nm{trg\_autorizacion\_bu} (C07))} & \parencite[pp.~15, 21, 23, 26 y~27]{BNA2026Tarjetas} & \nm{pruebas/\allowbreak postgres/\allowbreak pruebas.sql\allowbreak :336} & \textit{E:} consumido 0.00; disponible 15000.00 \newline \textit{O:} consumido 0.00; disponible 15000.00 & P & — \\
F12k & Una autorización solo se vincula al movimiento de su misma operación \textit{(Operación única (RN-BUS-002))} & \parencite[pp.~15, 21, 23, 26 y~27]{BNA2026Tarjetas} & \nm{pruebas/\allowbreak postgres/\allowbreak pruebas.sql\allowbreak :346} & \textit{E:} Rechazo: \nm{ck\_autorizacion\_3} \newline \textit{O:} Rechazo: new row for relation \textquotedbl{}autorizacion\textquotedbl{} violates check constraint \textquotedbl{}\nm{ck\_autorizacion\_3}\textquotedbl{} & P & — \\
F12l & Una autorización de ayer contabilizada hoy cuenta para el tope de ayer, no para el de hoy \textit{(Tope diario (RN-TOP-001))} & \parencites[pp.~15, 21, 23, 26 y~27]{BNA2026Tarjetas}[p.~4]{BNA2026Desembolso}; Decisión: el compromiso del tope pertenece al día de la autorización & \nm{pruebas/\allowbreak postgres/\allowbreak pruebas.sql\allowbreak :354} & \textit{E:} ayer 1000.00; hoy 1500.00; vinculada; contable 12500.00; disponible 12500.00 \newline \textit{O:} ayer 1000.00; hoy 1500.00; vinculada; contable 12500.00; disponible 12500.00 & P & — \\
F12m & El extorno hoy del movimiento de una autorización de ayer libera el tope de ayer, no el de hoy \textit{(Extorno (Tarjetas, C07))} & \parencites[pp.~15 y~28]{BNA2026Tarjetas}[p.~4]{BNA2026Desembolso}[pp.~16, 17, 25 y~32]{BNA2026Cuentas}{BNA2026Core}; Decisión: el compromiso del tope pertenece al día de la autorización & \nm{pruebas/\allowbreak postgres/\allowbreak pruebas.sql\allowbreak :373} & \textit{E:} ayer 0.00; hoy 400.00 \newline \textit{O:} ayer 0.00; hoy 400.00 & P & — \\
F12n & El extorno hoy de una autorización de ayer sin movimiento libera el tope de ayer y el saldo \textit{(\nm{trg\_autorizacion\_bu} (C07))} & \parencites[pp.~15, 21, 23, 26 y~27]{BNA2026Tarjetas}[p.~4]{BNA2026Desembolso}; Decisión: el compromiso del tope pertenece al día de la autorización & \nm{pruebas/\allowbreak postgres/\allowbreak pruebas.sql\allowbreak :388} & \textit{E:} ayer 0.00; hoy 400.00; disponible 14600.00 \newline \textit{O:} ayer 0.00; hoy 400.00; disponible 14600.00 & P & — \\
F13 & Una autorización no se borra; se extorna \textit{(\nm{trg\_autorizacion\_bd} (RN-ATM-003))} & \parencites{BNA2026Core}[p.~23]{BNA2026Tarjetas}[p.~13]{BNA2026Desembolso} & \nm{pruebas/\allowbreak postgres/\allowbreak pruebas.sql\allowbreak :405} & \textit{E:} Rechazo: ERR-TAR-009 \newline \textit{O:} Rechazo: ERR-TAR-009: tarjetas.autorizacion es de solo inserción (DELETE no permitido) & P & — \\
F14a & Si Clientes bloquea al cliente, se bloquean sus tarjetas \textit{(\nm{trg\_cliente\_ref\_au} (INT-T07))} & \parencites[pp.~23 y~31]{BNA2026Tarjetas}[p.~8]{BNA2026Clientes} & \nm{pruebas/\allowbreak postgres/\allowbreak pruebas.sql\allowbreak :411} & \textit{E:} BLOQUEADA por \nm{MODULO\_CLIENTES} \newline \textit{O:} BLOQUEADA por \nm{MODULO\_CLIENTES} & P & — \\
F14b & Si Clientes da de baja al cliente, se cancelan sus tarjetas \textit{(\nm{trg\_cliente\_ref\_au} (INT-T07))} & \parencites[pp.~23 y~31]{BNA2026Tarjetas}[p.~8]{BNA2026Clientes} & \nm{pruebas/\allowbreak postgres/\allowbreak pruebas.sql\allowbreak :420} & \textit{E:} CANCELADA \newline \textit{O:} CANCELADA & P & — \\
F15 & \nm{marcar\_tarjetas\_vencidas()} vence las tarjetas cuya fecha pasó \textit{(\nm{ev\_marcar\_tarjetas\_vencidas} (RN-TAR-003))} & \parencite[pp.~23 y~32]{BNA2026Tarjetas} & \nm{pruebas/\allowbreak postgres/\allowbreak pruebas.sql\allowbreak :425} & \textit{E:} 1 tarjeta; VENCIDA por SISTEMA \newline \textit{O:} 1 tarjeta; VENCIDA por SISTEMA & P & — \\
F16a & El crédito nace de una orden CONFIRMADA \textit{(Crédito (decisión I15))} & \parencites[p.~2]{BNA2026Cobranza}[p.~4]{BNA2026Solicitud}; Decisión: BNA-CREDITO; Decisión: decisión I15: el crédito nace de la orden confirmada & \nm{pruebas/\allowbreak postgres/\allowbreak pruebas.sql\allowbreak :439} & \textit{E:} Rechazo: BNA-CREDITO.*CONFIRMADA \newline \textit{O:} Rechazo: BNA-CREDITO: el crédito nace de una orden CONFIRMADA (la orden está PREPARADA) & P & — \\
F16b & El calendario del crédito es de la solicitud de su orden \textit{(Crédito (decisión I15))} & \parencites[p.~2]{BNA2026Cobranza}[p.~4]{BNA2026Solicitud}; Decisión: BNA-CREDITO; Decisión: decisión I15: el crédito nace de la orden confirmada & \nm{pruebas/\allowbreak postgres/\allowbreak pruebas.sql\allowbreak :442} & \textit{E:} Rechazo: BNA-CREDITO.*no pertenece \newline \textit{O:} Rechazo: BNA-CREDITO: el calendario d0000000-0000-4000-8000-000000000002 no pertenece a la solicitud b0000000-0000-4000-8000-000000000001 de la orden & P & — \\
F16c & El crédito usa el calendario OFICIAL \textit{(Crédito (Solicitud))} & \parencites[p.~2]{BNA2026Cobranza}[p.~4]{BNA2026Solicitud}; Decisión: BNA-CREDITO; Decisión: el crédito usa el calendario OFICIAL (supuesto S13) & \nm{pruebas/\allowbreak postgres/\allowbreak pruebas.sql\allowbreak :445} & \textit{E:} Rechazo: BNA-CREDITO.*OFICIAL \newline \textit{O:} Rechazo: BNA-CREDITO: el crédito usa el calendario OFICIAL, no el PRELIMINAR & P & — \\
F17a & Un pago no se aplica a cuotas de otro crédito \textit{(Aplicación de pagos)} & \parencite[pp.~3 y~8]{BNA2026Cobranza}; Decisión: BNA-PAGO & \nm{pruebas/\allowbreak postgres/\allowbreak pruebas.sql\allowbreak :448} & \textit{E:} Rechazo: BNA-PAGO.*no pertenece \newline \textit{O:} Rechazo: BNA-PAGO: la cuota no pertenece al calendario del crédito CR-0001 del pago & P & — \\
F17b & Solo se aplica un pago CONFIRMADO \textit{(Aplicación de pagos)} & \parencite[pp.~3 y~8]{BNA2026Cobranza}; Decisión: BNA-PAGO & \nm{pruebas/\allowbreak postgres/\allowbreak pruebas.sql\allowbreak :451} & \textit{E:} Rechazo: BNA-PAGO.*CONFIRMADO \newline \textit{O:} Rechazo: BNA-PAGO: solo se aplica un pago CONFIRMADO (el pago está PENDIENTE) & P & — \\
F18 & \nm{aplicar\_pago()} reparte el pago de la cuota más antigua a la siguiente \textit{(\nm{sp\_aplicar\_pago} (Cobranza))} & \parencite[p.~3]{BNA2026Cobranza}; Decisión: el pago se aplica de la cuota más antigua a la siguiente & \nm{pruebas/\allowbreak postgres/\allowbreak pruebas.sql\allowbreak :456} & \textit{E:} PAGADA 0.00; PAGADA 0.00; PARCIAL 3200.00 \newline \textit{O:} PAGADA 0.00; PAGADA 0.00; PARCIAL 3200.00 & P & — \\
F19 & Pagadas todas las cuotas, el crédito queda CANCELADO \textit{(Estado del crédito (Cobranza))} & \parencite[pp.~3 y~8]{BNA2026Cobranza} & \nm{pruebas/\allowbreak postgres/\allowbreak pruebas.sql\allowbreak :465} & \textit{E:} CANCELADO \newline \textit{O:} CANCELADO & P & — \\
F20 & Una aplicación de pago no se modifica \textit{(Aplicación de pagos)} & \parencites{BNA2026Core}[p.~23]{BNA2026Tarjetas}[p.~13]{BNA2026Desembolso}; Decisión: BNA-PAGO, BNA-SOLO-INSERCION; Decisión: una aplicación de pago no se modifica & \nm{pruebas/\allowbreak postgres/\allowbreak pruebas.sql\allowbreak :472} & \textit{E:} Rechazo: BNA-SOLO-INSERCION|BNA-PAGO \newline \textit{O:} Rechazo: BNA-PAGO: \nm{creditos.\allowbreak aplicacion\_pago} es de solo inserción (UPDATE no permitido) & P & — \\
F21a & Una cuenta BLOQUEADA no admite débitos \textit{(Estado de la cuenta (Cuentas))} & \parencites[pp.~16, 17, 25 y~32]{BNA2026Cuentas}{BNA2026Core}[pp.~15 y~28]{BNA2026Tarjetas}; Decisión: BNA-ESTADO-CUENTA & \nm{pruebas/\allowbreak postgres/\allowbreak pruebas.sql\allowbreak :477} & \textit{E:} Rechazo: BNA-ESTADO-CUENTA \newline \textit{O:} Rechazo: BNA-ESTADO-CUENTA: la cuenta CTA-PEN-2 (estado BLOQUEADA) no admite débitos & P & — \\
F21b & Una cuenta \nm{BLOQUEADA\_PREVENTIVA} admite créditos \textit{(Estado de la cuenta (Cuentas))} & \parencites[pp.~16, 17, 25 y~32]{BNA2026Cuentas}{BNA2026Core}[pp.~15 y~28]{BNA2026Tarjetas}; Decisión: una cuenta \nm{BLOQUEADA\_PREVENTIVA} admite créditos (Cuentas, p. 25, solo inhabilita débitos) & \nm{pruebas/\allowbreak postgres/\allowbreak pruebas.sql\allowbreak :481} & \textit{E:} saldo 5100.00 \newline \textit{O:} saldo 5100.00 & P & — \\
F22 & El extorno devuelve el monto al tope del día \textit{(Extorno (Tarjetas, C07))} & \parencites[pp.~16, 17, 25 y~32]{BNA2026Cuentas}{BNA2026Core}[pp.~15 y~28]{BNA2026Tarjetas} & \nm{pruebas/\allowbreak postgres/\allowbreak pruebas.sql\allowbreak :487} & \textit{E:} consumido hoy 0.00 \newline \textit{O:} consumido hoy 0.00 & P & — \\
F23 & El importe de un movimiento no se modifica \textit{(Movimiento inmutable)} & \parencite{BNA2026Core}; Decisión: BNA-MOVIMIENTO & \nm{pruebas/\allowbreak postgres/\allowbreak pruebas.sql\allowbreak :497} & \textit{E:} Rechazo: BNA-MOVIMIENTO \newline \textit{O:} Rechazo: BNA-MOVIMIENTO: solo se admite el paso de APROBADO a EXTORNADO & P & — \\
F24 & Un asiento sin líneas se rechaza al confirmar \textit{(Partida doble)} & \parencites{BNA2026Core}[p.~16]{BNA2026Cuentas}; Decisión: BNA-PARTIDA-DOBLE & \nm{pruebas/\allowbreak postgres/\allowbreak pruebas.sql\allowbreak :500} & \textit{E:} Rechazo: BNA-PARTIDA-DOBLE \newline \textit{O:} Rechazo: BNA-PARTIDA-DOBLE: el asiento 9a000000-0000-4000-8000-000000000003 no cuadra (debe 0.00, haber 0.00) & P & — \\
F25 & La vista enmascara el número y muestra el vencimiento MM/AA \textit{(\nm{v\_tarjeta\_resumen} (PCI-DSS))} & \parencite[pp.~23, 24, 27 y~28]{BNA2026Tarjetas} & \nm{pruebas/\allowbreak postgres/\allowbreak pruebas.sql\allowbreak :505} & \textit{E:} 4557 88** **** 1234; 01/31; operativa \newline \textit{O:} 4557 88** **** 1234; 01/31; operativa & P & — \\
F26 & La vista suma lo autorizado por tarjeta, canal, moneda y día \textit{(\nm{v\_consumo\_diario} (RN-TOP-001))} & \parencite[pp.~23, 24, 27 y~28]{BNA2026Tarjetas} & \nm{pruebas/\allowbreak postgres/\allowbreak pruebas.sql\allowbreak :512} & \textit{E:} ATM PEN hoy: 500.00 en 2 operaciones \newline \textit{O:} ATM PEN hoy: 500.00 en 2 operaciones & P & — \\
F27 & La vista muestra los intentos acumulados y restantes \textit{(\nm{v\_credenciales\_en\_riesgo} (RN-TAR-004))} & \parencite[pp.~23, 24, 27 y~28]{BNA2026Tarjetas} & \nm{pruebas/\allowbreak postgres/\allowbreak pruebas.sql\allowbreak :522} & \textit{E:} \nm{PIN\_ATM}: 1 fallido, 2 restantes \newline \textit{O:} \nm{PIN\_ATM}: 1 fallido, 2 restantes & P & — \\
F28 & La vista lista las órdenes con resultado incierto en el CORE \textit{(\nm{des\_v\_ordenes\_inciertas})} & \parencites[p.~12]{BNA2026Desembolso}[línea~92]{BNA2026Stack} & \nm{pruebas/\allowbreak postgres/\allowbreak pruebas.sql\allowbreak :529} & \textit{E:} DES-2026-000003 \newline \textit{O:} DES-2026-000003 (\nm{POR\_VERIFICAR}) & P & — \\
F29 & La vista lista las entregas que superan dos reintentos \textit{(\nm{des\_v\_entregas\_pendientes})} & \parencites[p.~12]{BNA2026Desembolso}[línea~92]{BNA2026Stack} & \nm{pruebas/\allowbreak postgres/\allowbreak pruebas.sql\allowbreak :534} & \textit{E:} 1 entrega, 3 intentos \newline \textit{O:} 1 entrega, 3 intentos & P & — \\
F30 & La vista cuadra lo ordenado con lo abonado por fecha y moneda \textit{(\nm{des\_v\_conciliacion\_diaria})} & \parencites[p.~12]{BNA2026Desembolso}[línea~92]{BNA2026Stack} & \nm{pruebas/\allowbreak postgres/\allowbreak pruebas.sql\allowbreak :539} & \textit{E:} PEN: ordenado 10000.00, abonado 10000.00, sin abono 0 \newline \textit{O:} PEN: ordenado 10000.00, abonado 10000.00, sin abono 0 & P & — \\
F31 & Existen los índices de consulta y ninguno repite a otro de clave foránea \textit{(Consultas de Tarjetas y Desembolso)} & \parencites[p.~30]{BNA2026Tarjetas}[p.~12]{BNA2026Desembolso} & \nm{pruebas/\allowbreak postgres/\allowbreak pruebas.sql\allowbreak :546} & \textit{E:} 8 índices de consulta; 0 redundantes \newline \textit{O:} 8 índices de consulta; 0 redundantes & P & — \\
R1 & Evento de auditoría sobre un registro no registrado como auditable \textit{(Pruebas negativas del DDL)} & Decisión: la entidad auditada es un catálogo (en Cobranza, p. 4, la referencia es libre) & \nm{pruebas/\allowbreak postgres/\allowbreak pruebas.sql\allowbreak :566} & \textit{E:} Rechazo: \nm{fk\_evento\_auditoria\_2} \newline \textit{O:} Rechazo: insert or update on table \textquotedbl{}\nm{evento\_auditoria}\textquotedbl{} violates foreign key constraint \textquotedbl{}\nm{fk\_evento\_auditoria\_2}\textquotedbl{} & P & — \\
R2 & Intento con coordenada para una credencial que no es la matriz \textit{(Pruebas negativas del DDL)} & \parencite[p.~20]{BNA2026Tarjetas} & \nm{pruebas/\allowbreak postgres/\allowbreak pruebas.sql\allowbreak :570} & \textit{E:} Rechazo: \nm{ck\_intento\_autenticacion\_4} \newline \textit{O:} Rechazo: new row for relation \textquotedbl{}\nm{intento\_autenticacion}\textquotedbl{} violates check constraint \textquotedbl{}\nm{ck\_intento\_autenticacion\_4}\textquotedbl{} & P & — \\
R3 & Intento con terminal y canal a la vez \textit{(Pruebas negativas del DDL)} & Decisión: arco exclusivo entre terminal y canal (Tarjetas, p. 20, admite ambos) & \nm{pruebas/\allowbreak postgres/\allowbreak pruebas.sql\allowbreak :574} & \textit{E:} Rechazo: \nm{ck\_intento\_autenticacion\_5} \newline \textit{O:} Rechazo: new row for relation \textquotedbl{}\nm{intento\_autenticacion}\textquotedbl{} violates check constraint \textquotedbl{}\nm{ck\_intento\_autenticacion\_5}\textquotedbl{} & P & — \\
R4 & Credencial de tipo \nm{TOKEN\_MATRIZ} \textit{(Pruebas negativas del DDL)} & \parencite[p.~18]{BNA2026Tarjetas} & \nm{pruebas/\allowbreak postgres/\allowbreak pruebas.sql\allowbreak :578} & \textit{E:} Rechazo: \nm{ck\_credencial\_1} \newline \textit{O:} Rechazo: new row for relation \textquotedbl{}credencial\textquotedbl{} violates check constraint \textquotedbl{}\nm{ck\_credencial\_1}\textquotedbl{} & P & — \\
R5 & Autorización sobre una cuenta inexistente \textit{(Pruebas negativas del DDL)} & \parencite[p.~21]{BNA2026Tarjetas} & \nm{pruebas/\allowbreak postgres/\allowbreak pruebas.sql\allowbreak :582} & \textit{E:} Rechazo: violates foreign key constraint \newline \textit{O:} Rechazo: insert or update on table \textquotedbl{}autorizacion\textquotedbl{} violates foreign key constraint \textquotedbl{}\nm{fk\_autorizacion\_5}\textquotedbl{} & P & — \\
R6 & Segundo cliente con el mismo DNI \textit{(Pruebas negativas del DDL)} & \parencite[p.~9]{BNA2026Clientes} & \nm{pruebas/\allowbreak postgres/\allowbreak pruebas.sql\allowbreak :586} & \textit{E:} Rechazo: \nm{uk\_cliente\_1} \newline \textit{O:} Rechazo: duplicate key value violates unique constraint \textquotedbl{}\nm{uk\_cliente\_1}\textquotedbl{} & P & — \\
R7 & Segunda orden CONFIRMADA para la misma aprobación \textit{(Pruebas negativas del DDL)} & \parencite[p.~12]{BNA2026Desembolso}; Decisión: BNA-ORDEN-UNICA & \nm{pruebas/\allowbreak postgres/\allowbreak pruebas.sql\allowbreak :592} & \textit{E:} Rechazo: BNA-ORDEN-UNICA|\nm{ux\_orden\_confirmada} \newline \textit{O:} Rechazo: BNA-ORDEN-UNICA: la solicitud b0000000-0000-4000-8000-000000000001 ya tiene una orden CONFIRMADA & P & — \\
R8 & Borrar una moneda en uso \textit{(Pruebas negativas del DDL)} & \parencite[línea~32]{BNA2026Stack}; Decisión: los catálogos en uso no se borran & \nm{pruebas/\allowbreak postgres/\allowbreak pruebas.sql\allowbreak :598} & \textit{E:} Rechazo: violates foreign key constraint \newline \textit{O:} Rechazo: update or delete on table \textquotedbl{}moneda\textquotedbl{} violates foreign key constraint \textquotedbl{}\nm{fk\_asiento\_contable\_1}\textquotedbl{} on table \textquotedbl{}\nm{asiento\_contable}\textquotedbl{} & P & — \\
M20b & Un asiento descuadrado se rechaza en un COMMIT real \textit{(Comprobación 20 (bug 7))} & \parencites{BNA2026Core}[p.~16]{BNA2026Cuentas}; MySQL: \nm{99\allowbreak :250}, \nm{99\allowbreak :224}; Decisión: BNA-PARTIDA-DOBLE & \nm{pruebas/\allowbreak postgres/\allowbreak run\_tests.sh\allowbreak :78} & \textit{E:} Rechazo: BNA-PARTIDA-DOBLE; asientos guardados: 0 \newline \textit{O:} Rechazo: BNA-PARTIDA-DOBLE: el asiento 9a000000-0000-4000-8000-000000000002 no cuadra (debe 100.00, haber 0.00); asientos guardados: 0 & P & — \\
C01 & Dos retiros simultáneos no dejan el saldo en negativo \textit{(Requisito 3 (E7))} & \parencites[pp.~16, 17, 25, 31 y~32]{BNA2026Cuentas}{BNA2026Core}[pp.~15 y~28]{BNA2026Tarjetas}; Decisión: BNA-SALDO & \nm{pruebas/\allowbreak postgres/\allowbreak run\_tests.sh\allowbreak :102} & \textit{E:} saldo 2000.00; segunda operación: BNA-SALDO \newline \textit{O:} saldo 2000.00; segunda operación: BNA-SALDO: saldo disponible insuficiente en CTA-PEN-2  & P & — \\
C02 & Dos retiros simultáneos en cajero no superan juntos el tope diario \textit{(RN-TOP-001 (bug 5))} & \parencites[pp.~16, 17, 25 y~32]{BNA2026Cuentas}{BNA2026Core}[pp.~15 y~28]{BNA2026Tarjetas} & \nm{pruebas/\allowbreak postgres/\allowbreak run\_tests.sh\allowbreak :112} & \textit{E:} consumido 1000.00; segunda operación: RN-TOP-001 \newline \textit{O:} consumido 1000.00; segunda operación: RN-TOP-001: tope diario de ATM excedido en CTA-PEN-1  & P & — \\
C03 & Dos autorizaciones simultáneas sin movimiento no superan juntas el tope diario \textit{(RN-TOP-001 (bug 5))} & \parencite[pp.~15, 21, 23, 26 y~27]{BNA2026Tarjetas} & \nm{pruebas/\allowbreak postgres/\allowbreak run\_tests.sh\allowbreak :122} & \textit{E:} autorizadas 1; segunda operación: ERR-TAR-007 \newline \textit{O:} autorizadas 1; segunda operación: ERR-TAR-007: tope diario excedido  & P & — \\
\addlinespace
\multicolumn{7}{@{}l}{\textit{Experimentos de comparación de los gestores}} \\
E1 & El DDL dentro de una transacción: MySQL lo confirma implícitamente; PostgreSQL lo deshace con ROLLBACK. & \parencite[sec.~15.3.3]{MySQL2024Manual} & \nm{experimentos/\allowbreak sgbd/\allowbreak experimentos.sh\allowbreak :71} & \textit{E:} MySQL: commit implícito; PG: se deshizo \newline \textit{O:} MySQL 8.4: la tabla sigue existiendo (commit implícito); PG 16 y 17: la tabla se deshizo & P & P \\
E2 (Clave foránea diferible) & MySQL no admite claves foráneas diferibles; PostgreSQL sí. & \parencites[sec.~15.1.20.5]{MySQL2024Manual}[CREATE TABLE]{PostgreSQL2023Docs} & \nm{experimentos/\allowbreak sgbd/\allowbreak experimentos.sh\allowbreak :82} & \textit{E:} MySQL: error in your SQL syntax; PG: \textasciicircum{}aceptado\$ \newline \textit{O:} MySQL 8.4: You have an error in your SQL syntax; check the manual that corresponds to your MySQL serv; PG 16 y 17: aceptado & P & P \\
E2 (Índice único parcial) & MySQL no admite índices únicos parciales; PostgreSQL sí. & \parencite[sec.~11.8]{PostgreSQL2023Docs} & \nm{experimentos/\allowbreak sgbd/\allowbreak experimentos.sh\allowbreak :82} & \textit{E:} MySQL: error in your SQL syntax; PG: \textasciicircum{}aceptado\$ \newline \textit{O:} MySQL 8.4: You have an error in your SQL syntax; check the manual that corresponds to your MySQL serv; PG 16 y 17: aceptado & P & P \\
E2 (Dominio) & MySQL no admite dominios; PostgreSQL sí. & \parencite[sec.~8.18]{PostgreSQL2023Docs} & \nm{experimentos/\allowbreak sgbd/\allowbreak experimentos.sh\allowbreak :82} & \textit{E:} MySQL: error in your SQL syntax; PG: \textasciicircum{}aceptado\$ \newline \textit{O:} MySQL 8.4: You have an error in your SQL syntax; check the manual that corresponds to your MySQL serv; PG 16 y 17: aceptado & P & P \\
E3 & MySQL no puede declarar un trigger diferido hasta la confirmación; en PostgreSQL, un trigger de restricción diferido rechaza al confirmar el asiento descuadrado y guarda el que cuadra. & \parencite[CREATE TRIGGER]{PostgreSQL2023Docs} & \nm{experimentos/\allowbreak sgbd/\allowbreak experimentos.sh\allowbreak :99} & \textit{E:} MySQL: CREATE CONSTRAINT TRIGGER: .*error in your SQL syntax; PG: asiento 2 descuadrado.*asientos guardados: 1\$ \newline \textit{O:} MySQL 8.4: CREATE CONSTRAINT TRIGGER: You have an error in your SQL syntax; check the manual that corresponds to your MyS; PG 16 y 17: asiento 2 descuadrado asientos guardados: 1 & P & P \\
E4 & Con \nm{sql\_mode} = '', MySQL guarda valores inválidos ajustados, solo con advertencias; PostgreSQL los rechaza siempre. & \parencite[sec.~7.1.11]{MySQL2024Manual} & \nm{experimentos/\allowbreak sgbd/\allowbreak experimentos.sh\allowbreak :122} & \textit{E:} MySQL: guardado sin error: e='', d=999\textbackslash{}.99, f=0000-00-00; advertencias: [1-9]; PG: numeric field overflow.*out of range.*filas guardadas: 0 \newline \textit{O:} MySQL 8.4: guardado sin error: e='', d=999.99, f=0000-00-00; advertencias: 3; PG 16 y 17: numeric field overflow date/time field value out of range: \textquotedbl{}2026-02-31\textquotedbl{} filas guardadas: 0 & P & P \\
E5 (Marca de tiempo de 2039) & El TIMESTAMP de MySQL no llega a 2039; el TIMESTAMPTZ de PostgreSQL sí. & \parencites[sec.~13.2.2]{MySQL2024Manual}[sec.~8.5]{PostgreSQL2023Docs} & \nm{experimentos/\allowbreak sgbd/\allowbreak experimentos.sh\allowbreak :136} & \textit{E:} MySQL: Incorrect datetime value; PG: TIMESTAMPTZ guarda 2039 \newline \textit{O:} MySQL 8.4: TIMESTAMP: Incorrect datetime value: '2039-01-01 00:00:00' for column 't' at row 1; PG 16 y 17: TIMESTAMPTZ guarda 2039-01-01 05:00:00+00 & P & P \\
E5 (BOOLEAN = 5) & El BOOLEAN de MySQL es un TINYINT(1) y acepta 5; el de PostgreSQL lo rechaza. & \parencites[sec.~13.1.1]{MySQL2024Manual}[sec.~8.6]{PostgreSQL2023Docs} & \nm{experimentos/\allowbreak sgbd/\allowbreak experimentos.sh\allowbreak :136} & \textit{E:} MySQL: BOOLEAN guarda 5; PG: type boolean \newline \textit{O:} MySQL 8.4: BOOLEAN guarda 5; PG 16 y 17: BOOLEAN: column \textquotedbl{}x\textquotedbl{} is of type boolean but expression is of type integer & P & P \\
E5 (UUID 'c2') & MySQL no tiene tipo UUID y CHAR(36) acepta 'c2'; el UUID de PostgreSQL lo rechaza. & \parencite[sec.~8.12]{PostgreSQL2023Docs} & \nm{experimentos/\allowbreak sgbd/\allowbreak experimentos.sh\allowbreak :136} & \textit{E:} MySQL: CHAR\textbackslash{}(36\textbackslash{}) guarda 'c2'; PG: invalid input syntax for type uuid \newline \textit{O:} MySQL 8.4: CHAR(36) guarda 'c2'; PG 16 y 17: UUID: invalid input syntax for type uuid: \textquotedbl{}c2\textquotedbl{} & P & P \\
E6 & Un borrado en cascada no dispara el trigger de auditoría en MySQL y sí en PostgreSQL. & \parencites[sec.~15.1.20.5]{MySQL2024Manual}[CREATE TRIGGER]{PostgreSQL2023Docs} & \nm{experimentos/\allowbreak sgbd/\allowbreak experimentos.sh\allowbreak :155} & \textit{E:} MySQL: bitácora: 0\$; PG: bitácora: 1\$ \newline \textit{O:} MySQL 8.4: hijos borrados: 1; filas en la bitácora: 0; PG 16 y 17: hijos borrados: 1; filas en la bitácora: 1 & P & P \\
E7 (READ COMMITTED) & Con el aislamiento por omisión de PostgreSQL se pierde la actualización (saldo final 50). & \parencite[cap.~13]{PostgreSQL2023Docs} & \nm{experimentos/\allowbreak sgbd/\allowbreak experimentos.sh\allowbreak :178} & \textit{E:} PG: A confirma; saldo final 50 \newline \textit{O:} PG 16 y 17: READ COMMITTED: A confirma; saldo final 50.00 & P & — \\
E7 (REPEATABLE READ) & Con el aislamiento por omisión de MySQL se pierde la actualización; en PostgreSQL, REPEATABLE READ aborta la segunda transacción, que debe reintentarse. & \parencites[sec.~17.7.2.1]{MySQL2024Manual}[cap.~13]{PostgreSQL2023Docs} & \nm{experimentos/\allowbreak sgbd/\allowbreak experimentos.sh\allowbreak :178} & \textit{E:} MySQL: A confirma; saldo final 50; PG: could not serialize \newline \textit{O:} MySQL 8.4: REPEATABLE READ: A confirma; saldo final 50.00; PG 16 y 17: REPEATABLE READ: could not serialize access due to concurrent update; saldo final 70.00 & P & P \\
E7 (SERIALIZABLE) & Con SERIALIZABLE, MySQL hace esperar a la segunda transacción (saldo final 20) y PostgreSQL la aborta para que se reintente. & \parencites[sec.~17.7.2.1]{MySQL2024Manual}[cap.~13]{PostgreSQL2023Docs} & \nm{experimentos/\allowbreak sgbd/\allowbreak experimentos.sh\allowbreak :178} & \textit{E:} MySQL: A confirma; saldo final 20; PG: could not serialize \newline \textit{O:} MySQL 8.4: SERIALIZABLE: A confirma; saldo final 20.00; PG 16 y 17: SERIALIZABLE: could not serialize access due to concurrent update; saldo final 70.00 & P & P \\
\addlinespace
\multicolumn{7}{@{}l}{\textit{Comprobaciones sobre la línea base MySQL (commit ab6a6ec)}} \\
MY01 & El script de carga encuentra el archivo SQL \textit{(Comprobación 1 (bug 1))} & MySQL: \nm{cargar\_bna.ps1\allowbreak :11}, \nm{GUIA\_EJECUCION.md\allowbreak :83}, \nm{README.md\allowbreak :53} & \nm{repro\_mysql/\allowbreak run.sh\allowbreak :106} & \textit{E:} Que apunte a \nm{Trabajo\_hecho}\textbackslash{}BD\textbackslash{}\nm{99\_bna\_unificado.\allowbreak sql} \newline \textit{O:} MySQL 8.4: Busca \nm{Trabajo\_hecho}\textbackslash{}\nm{99\_bna\_unificado.\allowbreak sql}, que no existe (el archivo esta en BD\textbackslash{}) & — & F \\
MY02 & El script unificado carga sin errores en MySQL 8.4 \textit{(Comprobación 2)} & MySQL: \nm{99\allowbreak :23}, \nm{99\allowbreak :29} & \nm{repro\_mysql/\allowbreak run.sh\allowbreak :115} & \textit{E:} Termina sin ningun ERROR \newline \textit{O:} MySQL 8.4: Termino sin errores (22 avisos de obsolescencia, inofensivos: TINYINT(1) y \nm{sql\_mode}) & — & OK \\
MY03 & La base tiene los objetos que promete la guia \textit{(Comprobación 3)} & MySQL: \nm{GUIA\_EJECUCION.md\allowbreak :9}, \nm{verificar\_bna.ps1\allowbreak :33} & \nm{repro\_mysql/\allowbreak sql/\allowbreak 00\_conteos.sql\allowbreak :20} & \textit{E:} 66 / 9 / 7 / 3 / 1, como dicen \nm{GUIA\_EJECUCION.\allowbreak md} y \nm{verificar\_bna.\allowbreak ps1} \newline \textit{O:} MySQL 8.4: 66 / 9 / 7 / 3 / 1 & — & OK \\
MY04 & Un deposito actualiza el saldo de la cuenta \textit{(Comprobación 4 (bug 3))} & MySQL: \nm{99\allowbreak :1996}, \nm{99\allowbreak :388}, \nm{99\allowbreak :1811} & \nm{repro\_mysql/\allowbreak sql/\allowbreak 01\_estado\_demo.sql\allowbreak :9} & \textit{E:} \nm{saldo\_contable} = 500.00 \newline \textit{O:} MySQL 8.4: \nm{saldo\_contable} = 0.00 & — & F \\
MY05 & Pagar capital reduce la deuda del credito \textit{(Comprobación 5 (bug 4))} & MySQL: \nm{99\allowbreak :1865}, \nm{99\allowbreak :1244} & \nm{repro\_mysql/\allowbreak sql/\allowbreak 01\_estado\_demo.sql\allowbreak :20} & \textit{E:} \nm{saldo\_capital} = 9700.00 \newline \textit{O:} MySQL 8.4: \nm{saldo\_capital} = 10000.00 & — & F \\
MY06 & \nm{fn\_uuid\_v7()} guarda la hora correcta \textit{(Comprobación 6 (bug 2))} & MySQL: \nm{99\allowbreak :1699}, \nm{99\allowbreak :1700} & \nm{repro\_mysql/\allowbreak sql/\allowbreak 01\_estado\_demo.sql\allowbreak :36} & \textit{E:} La hora del UUID coincide con la hora actual (diferencia 0 ms, a lo sumo 1-2 ms) \newline \textit{O:} MySQL 8.4: La hora del UUID esta 591 ms en el futuro & — & F \\
MY07 & Los UUID v7 se ordenan por fecha de creacion \textit{(Comprobación 7 (bug 2))} & MySQL: \nm{99\allowbreak :1700}, \nm{99\allowbreak :18}, \nm{diccionario\_datos.md\allowbreak :7} & \nm{repro\_mysql/\allowbreak sql/\allowbreak 01\_estado\_demo.sql\allowbreak :50} & \textit{E:} El UUID de 12:00:00.900 se ordena antes que el de 12:00:01.100 \newline \textit{O:} MySQL 8.4: El UUID de 12:00:00.900 se ordena DESPUES que el de 12:00:01.100 & — & F \\
MY08 & Los IDs tienen que ser UUID validos \textit{(Comprobación 8 (bug 7))} & MySQL: \nm{99\allowbreak :289}, \nm{99\allowbreak :6} & \nm{repro\_mysql/\allowbreak sql/\allowbreak 02\_integridad.sql\allowbreak :13} & \textit{E:} Rechazo: un UUID tiene 36 caracteres con formato 8-4-4-4-12 \newline \textit{O:} MySQL 8.4: MySQL lo acepto sin error & — & F \\
MY09 & Una orden de desembolso respeta su solicitud \textit{(Comprobación 9 (bug 7))} & MySQL: \nm{99\allowbreak :1106}, \nm{99\allowbreak :1110} & \nm{repro\_mysql/\allowbreak sql/\allowbreak 02\_integridad.sql\allowbreak :24} & \textit{E:} Rechazo: el cliente, la moneda y el monto salen de la solicitud aprobada \newline \textit{O:} MySQL 8.4: MySQL lo acepto sin error & — & F \\
MY10 & Un movimiento de cuenta es coherente \textit{(Comprobación 10 (bug 7))} & MySQL: \nm{99\allowbreak :450}, \nm{99\allowbreak :430}, \nm{99\allowbreak :433} & \nm{repro\_mysql/\allowbreak sql/\allowbreak 02\_integridad.sql\allowbreak :36} & \textit{E:} Rechazo: un retiro es siempre debito, la moneda es la de la cuenta y el saldo no puede ser negativo \newline \textit{O:} MySQL 8.4: MySQL lo acepto sin error & — & F \\
MY11 & El saldo disponible no supera al contable \textit{(Comprobación 11 (bug 7))} & MySQL: \nm{99\allowbreak :401} & \nm{repro\_mysql/\allowbreak sql/\allowbreak 02\_integridad.sql\allowbreak :47} & \textit{E:} Rechazo: no se puede disponer de mas dinero del que hay (CHECK del informe) \newline \textit{O:} MySQL 8.4: MySQL lo acepto sin error & — & F \\
MY12 & El total de una cuota es la suma de sus partes \textit{(Comprobación 12 (bug 7))} & MySQL: \nm{99\allowbreak :1269}, \nm{99\allowbreak :1277} & \nm{repro\_mysql/\allowbreak sql/\allowbreak 02\_integridad.sql\allowbreak :59} & \textit{E:} Rechazo, o mejor aun: que \nm{cuota\_total} no se guarde y se calcule (1 + 1 = 2) \newline \textit{O:} MySQL 8.4: MySQL lo acepto sin error & — & F \\
MY13 & \nm{sp\_aplicar\_pago} no aplica mas de lo pagado ni cuenta la mora como cuota \textit{(Comprobación 13 (bug 4))} & MySQL: \nm{99\allowbreak :1853}, \nm{99\allowbreak :1867} & \nm{repro\_mysql/\allowbreak sql/\allowbreak 02\_integridad.sql\allowbreak :70} & \textit{E:} Error: no se pueden aplicar 9999 de un pago de 1; y pagar mora no cancela capital ni interes \newline \textit{O:} MySQL 8.4: Se aplico sin error; la cuota quedo PAGADA con \nm{monto\_pagado} = 9999.00 & — & F \\
MY14 & Una solicitud genera un solo credito \textit{(Comprobación 14 (bug 7))} & MySQL: \nm{99\allowbreak :1251}, \nm{99\allowbreak :1238} & \nm{repro\_mysql/\allowbreak sql/\allowbreak 02\_integridad.sql\allowbreak :85} & \textit{E:} Rechazo: en el informe el credito es 1:1 con su orden y su calendario \newline \textit{O:} MySQL 8.4: MySQL lo acepto sin error & — & F \\
MY15 & Una tarjeta tiene datos validos \textit{(Comprobación 15 (bug 7))} & MySQL: \nm{99\allowbreak :527}, \nm{99\allowbreak :531} & \nm{repro\_mysql/\allowbreak sql/\allowbreak 02\_integridad.sql\allowbreak :96} & \textit{E:} Rechazo: ultimos4 son 4 digitos y el vencimiento es posterior a la emision \newline \textit{O:} MySQL 8.4: MySQL lo acepto sin error & — & F \\
MY16 & Una tarjeta matriz cuelga de una credencial de tipo matriz \textit{(Comprobación 16 (bug 7))} & MySQL: \nm{99\allowbreak :563}, \nm{99\allowbreak :587} & \nm{repro\_mysql/\allowbreak sql/\allowbreak 02\_integridad.sql\allowbreak :108} & \textit{E:} Rechazo: el PIN tiene 4 digitos y la matriz no pertenece a un PIN \newline \textit{O:} MySQL 8.4: MySQL lo acepto sin error & — & F \\
MY17 & La coordenada pedida existe en la matriz \textit{(Comprobación 17 (bug 7))} & MySQL: \nm{99\allowbreak :639}, \nm{99\allowbreak :514} & \nm{repro\_mysql/\allowbreak sql/\allowbreak 02\_integridad.sql\allowbreak :118} & \textit{E:} Rechazo: fila A-H y columna 1-8 \newline \textit{O:} MySQL 8.4: MySQL lo acepto sin error & — & F \\
MY18 & Un caso de captacion apunta a una solicitud real \textit{(Comprobación 18 (bug 7))} & MySQL: \nm{99\allowbreak :766}, \nm{99\allowbreak :881} & \nm{repro\_mysql/\allowbreak sql/\allowbreak 02\_integridad.sql\allowbreak :129} & \textit{E:} Rechazo por clave foranea \newline \textit{O:} MySQL 8.4: MySQL lo acepto sin error & — & F \\
MY19 & La auditoria solo registra eventos sobre registros reales \textit{(Comprobación 19)} & MySQL: \nm{99\allowbreak :259}, \nm{99\allowbreak :260} & \nm{repro\_mysql/\allowbreak sql/\allowbreak 02\_integridad.sql\allowbreak :140} & \textit{E:} Rechazo: en el informe la entidad auditada es una FK a \nm{ENTIDAD\_AUDITADA} \newline \textit{O:} MySQL 8.4: MySQL lo acepto sin error & — & F \\
MY20 & Un asiento contable cuadra y cae en su periodo \textit{(Comprobación 20 (bug 7))} & MySQL: \nm{99\allowbreak :250}, \nm{99\allowbreak :224} & \nm{repro\_mysql/\allowbreak sql/\allowbreak 02\_integridad.sql\allowbreak :158} & \textit{E:} Rechazo: en partida doble el debe y el haber suman lo mismo, y la fecha cae dentro del periodo \newline \textit{O:} MySQL 8.4: MySQL lo acepto sin error & — & F \\
MY21 & Una reversion la autoriza otra persona \textit{(Comprobación 21)} & MySQL: \nm{99\allowbreak :1201}, \nm{99\allowbreak :1496}, \nm{99\allowbreak :1497} & \nm{repro\_mysql/\allowbreak sql/\allowbreak 02\_integridad.sql\allowbreak :172} & \textit{E:} Rechazo: quien solicita y quien autoriza deben ser personas distintas \newline \textit{O:} MySQL 8.4: MySQL lo acepto sin error & — & F \\
MY22 & Un retiro se cuenta una sola vez en el tope diario \textit{(Comprobación 22 (bug 5))} & MySQL: \nm{99\allowbreak :1801}, \nm{99\allowbreak :1816} & \nm{repro\_mysql/\allowbreak sql/\allowbreak 03\_topes\_y\_bloqueo.sql\allowbreak :16} & \textit{E:} Acumulado del dia = 600.00 \newline \textit{O:} MySQL 8.4: Acumulado del dia = 1200.00 & — & F \\
MY23 & Una autorizacion aprobada despues tambien cuenta para el tope \textit{(Comprobación 23 (bug 5))} & MySQL: \nm{99\allowbreak :1796}, \nm{99\allowbreak :1799} & \nm{repro\_mysql/\allowbreak sql/\allowbreak 03\_topes\_y\_bloqueo.sql\allowbreak :31} & \textit{E:} El acumulado sube de 1200.00 a 6200.00 \newline \textit{O:} MySQL 8.4: El acumulado sigue en 1200.00 & — & F \\
MY24 & Se respeta el tope diario del cajero \textit{(Comprobación 24 (bug 5))} & MySQL: \nm{99\allowbreak :1802}, \nm{99\allowbreak :1412} & \nm{repro\_mysql/\allowbreak sql/\allowbreak 03\_topes\_y\_bloqueo.sql\allowbreak :38} & \textit{E:} Rechazo: supera el tope diario del canal \newline \textit{O:} MySQL 8.4: MySQL lo acepto sin error & — & F \\
MY25 & \nm{v\_tope\_restante\_hoy} muestra el tope de cada cuenta \textit{(Comprobación 25 (bug 5))} & MySQL: \nm{99\allowbreak :1568}, \nm{99\allowbreak :1569} & \nm{repro\_mysql/\allowbreak sql/\allowbreak 03\_topes\_y\_bloqueo.sql\allowbreak :46} & \textit{E:} Una fila por cuenta y canal, con su tope restante \newline \textit{O:} MySQL 8.4: Una sola fila para todo el canal ATM: consumido 5600.00 entre todas las cuentas & — & F \\
MY26 & Al bloquear una tarjeta se guarda el motivo \textit{(Comprobación 26 (bug 6))} & MySQL: \nm{99\allowbreak :1827}, \nm{99\allowbreak :1832}, \nm{99\allowbreak :1736} & \nm{repro\_mysql/\allowbreak sql/\allowbreak 03\_topes\_y\_bloqueo.sql\allowbreak :62} & \textit{E:} Historial: motivo = Robo reportado, origen = CLIENTE \newline \textit{O:} MySQL 8.4: Historial: motivo = Cambio de estado, origen = SISTEMA & — & F \\
MY27 & El diccionario de datos describe la base real \textit{(Comprobación 27)} & MySQL: \nm{diccionario\_datos.md\allowbreak :10} & \nm{repro\_mysql/\allowbreak run.sh\allowbreak :145} & \textit{E:} Las mismas columnas en las mismas tablas \newline \textit{O:} MySQL 8.4: Coinciden las 521 columnas & — & OK \\
MY28 & El diagrama Mermaid muestra todas las relaciones \textit{(Comprobación 28)} & MySQL: \nm{mer\_bna.mermaid\allowbreak :2}, \nm{mer\_bna.mermaid\allowbreak :63} & \nm{repro\_mysql/\allowbreak run.sh\allowbreak :160} & \textit{E:} Una relacion por clave foranea (109) y todas las tablas presentes \newline \textit{O:} MySQL 8.4: 68 relaciones de 109; faltan las tablas: \nm{auditoria\_evento}, \nm{eva\_regla} & — & F \\
